\pagebreak

\section{Theorems}\label{sec:theorems}

\subsection{Fisher–Tippett–Gnedenko Theorem}\label{subsec:theorem_gev}
For completeness, we state below the Fisher–Tippett–Gnedenko theorem, which forms the foundation of block-maxima approaches in EVT. 
We omit the proof, which can be found in standard references such as 
\citet{gnedenko1943distribution} and \citet{Coles2001}.


\begin{theorem}[Fisher–Tippett–Gnedenko] \label{thm:gev}
Let \( X_1, \ldots, X_n \sim F \) i.i.d., and define \( M_n = \max(X_1, \ldots, X_n) \). 
If sequences \( {a_n} > 0 \), \( {b_n} \in \mathbb{R} \) exist such that
\begin{equation}
\frac{M_n - b_n}{a_n} \xrightarrow{d} G_\gamma(x),
\label{eqn:gev-limit}
\end{equation}
then the limiting distribution \( G_\gamma(x) \) belongs to the \textbf{Generalised Extreme Value (GEV)} family:
\begin{equation}
G_\gamma(x) = \exp\left\{ -\left(1 + \gamma \frac{x - \mu}{\sigma} \right)^{-1/\gamma} \right\}, 
\quad 1 + \gamma \frac{x - \mu}{\sigma} > 0,
\label{eqn:gev}
\end{equation}
where \( \mu \in \mathbb{R} \), \( \sigma > 0 \) and \( \gamma \in \mathbb{R} \) 
are location, scale and shape parameters, respectively.
\end{theorem}
